\documentclass[../numerical_methods.tex]{subfiles}
\begin{document}
\section{Aula 13 - ?? de Maio, 2026}
\subsection{Motivações}
\begin{itemize}
	\item Relacionando métodos numéricos.
\end{itemize}
\subsection{Eliminação de Gauss e Decomposição LU}
Nesta aula, iremos mostrar que a Eliminação de Gauss e a Decomposição LU são, na verdade, métodos numéricos equivalentes (em particular, possuem a mesma complexidade de \(\frac{2}{3}n^{3}\) flops.
Para isso, veremos a versão de matrizes elementares da eliminação de Gauss: dada uma equação
\[
	Ax = b, \quad A = \begin{bmatrix}
		a_{11}^{(1)} & a_{12}^{(1)} & a_{13}^{(1)} & \dotsc & a_{1n}^{(1)} \\
		a_{21}^{(1)} & a_{22}^{(1)} & a_{23}^{(1)} & \dotsc & a_{2n}^{(1)} \\
		a_{31}^{(1)} & a_{32}^{(1)} & a_{33}^{(1)} & \dotsc & a_{3n}^{(1)} \\
		\vdots       & \vdots       & \vdots       & \ddots & \vdots       \\
		a_{n1}^{(1)} & a_{n2}^{(1)} & a_{n3}^{(1)} & \dotsc & a_{nn}^{(1)} \\
	\end{bmatrix}
\]
começamos multiplicando ambos os lados por uma matriz \(M_{1}\), mantendo o sistema equivalente:
\[
	M_1 Ax = M_1 b,
\]
construída com o objetivo de ``matricializar'' a operação de anular os termos abaixo do pivô. Assim, queremos
\[
	M_{1}A = A_{1} = \begin{bmatrix}
		a_{11}^{(1)} & a_{12}^{(1)} & a_{13}^{(1)} & \dotsc & a_{1n}^{(1)} \\
		0            & a_{22}^{(2)} & a_{23}^{(2)} & \dotsc & a_{2n}^{(2)} \\
		0            & a_{32}^{(2)} & a_{33}^{(2)} & \dotsc & a_{3n}^{(2)} \\
		\vdots       & \vdots       & \vdots       & \ddots & \vdots       \\
		0            & a_{n2}^{(2)} & a_{n3}^{(2)} & \dotsc & a_{nn}^{(2)} \\
	\end{bmatrix}.
\]
Logo, \(M_{1}\) precisa ter a forma
\[
	M_{1} = \begin{bmatrix}
		1      &        &        &        &   \\
		m_{21} & 1      &        &        &   \\
		m_{31} &        & 1      &        &   \\
		\vdots & \vdots & \vdots & \ddots &   \\
		m_{n1} &        &        &        & n
	\end{bmatrix}
\]
\end{document}
